% =============================================================
%  Void Linux + Sway — A Comprehensive Guide
%  Compile with: pdflatex void-sway-guide-en.tex
% =============================================================
\documentclass[12pt, a4paper]{book}

\usepackage[utf8]{inputenc}
\usepackage[T1]{fontenc}
\usepackage[english]{babel}

% Page layout
\usepackage[a4paper, top=2.5cm, bottom=2.5cm, left=3cm, right=3cm]{geometry}

% Code listings
\usepackage{listings}
\usepackage{xcolor}
\lstset{
  basicstyle=\ttfamily\small,
  backgroundcolor=\color{gray!10},
  frame=single,
  breaklines=true,
  breakatwhitespace=true,
  showstringspaces=false,
  tabsize=2,
  xleftmargin=8pt,
  xrightmargin=8pt,
  framesep=4pt,
  commentstyle=\color{gray},
  keywordstyle=\bfseries,
}

% Tables
\usepackage{booktabs}
\usepackage{tabularx}
\usepackage{array}
\renewcommand{\arraystretch}{1.3}

% Misc
\usepackage{hyperref}
\usepackage{enumitem}
\usepackage{parskip}
%\usepackage{microtype}

\hypersetup{
  colorlinks=true,
  linkcolor=black,
  urlcolor=blue,
  pdftitle={Void Linux + Sway: A Comprehensive Guide},
}

% Simple box for notes/warnings using a framed environment
\usepackage{mdframed}
\newmdenv[linecolor=black, linewidth=0.6pt, innerleftmargin=8pt,
  innerrightmargin=8pt, innertopmargin=6pt, innerbottommargin=6pt,
  backgroundcolor=gray!8]{notebox}
\newmdenv[linecolor=black, linewidth=1pt, innerleftmargin=8pt,
  innerrightmargin=8pt, innertopmargin=6pt, innerbottommargin=6pt,
  backgroundcolor=gray!5]{warningbox}

% Column types for tables
\newcolumntype{L}{>{\raggedright\arraybackslash}X}

% =============================================================
\begin{document}

% Title page
\begin{titlepage}
  \centering
  \vspace*{4cm}
  {\Huge\bfseries Void Linux + Sway\par}
  \vspace{0.8cm}
  {\Large A Comprehensive Setup Guide\par}
  \vspace{1.5cm}
  \rule{0.5\textwidth}{0.5pt}
  \vspace{1.5cm}\par
  {\normalsize
  A personal reference guide for building the environment from scratch\\
  on Void Linux with the Sway window manager.\\[0.4em]
  Built from real-world experience and official documentation.\\[0.4em]
  Each chapter is self-contained --- read only what you need.
  \par}
  \vspace{2cm}
  {\normalsize\itshape Recommended installation order:\\[0.3em]
  01 $\to$ 02 $\to$ 03 $\to$ 04 $\to$ 05 $\to$ 06 $\to$ 07}
  \vfill
  \rule{0.4\textwidth}{0.3pt}
\end{titlepage}

\tableofcontents
\clearpage

% =============================================================
\chapter*{Void Linux Quick Reference}
\addcontentsline{toc}{chapter}{Void Linux Quick Reference}

\begin{lstlisting}
# Install a package
sudo xbps-install PACKAGE

# Enable a service
sudo ln -s /etc/sv/SERVICENAME /var/service/

# After any "usermod -aG" --> log out and back in

# Launch Sway correctly
dbus-run-session sway

# elogind vs seatd --> pick ONE, they cannot run together
\end{lstlisting}

% =============================================================
\chapter{Basic Installation}
% =============================================================

\section{Update the System First}

\begin{lstlisting}
sudo xbps-install -Su
\end{lstlisting}

\section{Essential Build Tools}

\begin{lstlisting}
sudo xbps-install base-devel git make ca-certificates wget curl
\end{lstlisting}

\section{GPU Drivers (Mesa without Xorg)}

\begin{notebox}
On Wayland you do \textbf{not} need \texttt{xorg-server} or \texttt{xf86-video-intel}.
\end{notebox}

\subsection{Intel (e.g.\ Dell Latitude E6430)}
\begin{lstlisting}
sudo xbps-install mesa-dri intel-video-accel \
  libva-utils mesa-vaapi mesa-vdpau \
  vulkan-loader mesa-vulkan-intel
\end{lstlisting}

\subsection{AMD}
\begin{lstlisting}
sudo xbps-install mesa-dri mesa-vaapi mesa-vdpau \
  vulkan-loader mesa-vulkan-radeon
\end{lstlisting}

\subsection{NVIDIA}
\begin{lstlisting}
sudo xbps-install nvidia nvidia-libs
\end{lstlisting}

\section{Session Manager (Choose One Only!)}

\begin{warningbox}
\textbf{Warning:} \texttt{elogind} and \texttt{seatd} cannot run together --- choose one.
\end{warningbox}

\subsection{Option A: elogind (recommended for laptops)}
\begin{lstlisting}
sudo xbps-install elogind
sudo ln -s /etc/sv/elogind /var/service/
\end{lstlisting}

\textbf{elogind advantages:}
\begin{itemize}[leftmargin=*]
  \item Creates \texttt{XDG\_RUNTIME\_DIR} automatically
  \item Supports power management (suspend, hibernate)
  \item More stable with most applications
\end{itemize}

\subsection{Option B: seatd (lighter and simpler)}
\begin{lstlisting}
sudo xbps-install seatd
sudo ln -s /etc/sv/seatd /var/service/
sudo usermod -aG _seatd $USER
\end{lstlisting}

\textbf{seatd advantages:}
\begin{itemize}[leftmargin=*]
  \item Much lighter
  \item Manages only the seat (not \texttt{XDG\_RUNTIME\_DIR})
  \item Requires \texttt{pam\_rundir} alongside it
\end{itemize}

\section{XDG\_RUNTIME\_DIR}

\subsection{If you chose elogind}
Nothing needed --- it is created automatically.

\subsection{If you chose seatd}
\begin{lstlisting}
sudo xbps-install pam_rundir
\end{lstlisting}
Add to \texttt{/etc/pam.d/system-login}:
\begin{lstlisting}
-session optional pam_rundir.so
\end{lstlisting}

\section{D-Bus}
\begin{lstlisting}
sudo xbps-install dbus
sudo ln -s /etc/sv/dbus /var/service/
\end{lstlisting}

\section{Default User Directories}
\begin{lstlisting}
sudo xbps-install xdg-user-dirs
xdg-user-dirs-update
\end{lstlisting}

\section{Sway and Core Tools}
\begin{lstlisting}
sudo xbps-install \
  sway swaybg swayidle swaylock \
  xorg-server-xwayland \
  wlroots wl-clipboard

# Terminal
sudo xbps-install tilix   # or foot (lighter and faster)
sudo xbps-install foot

# App Launcher
sudo xbps-install fuzzel   # or wofi
sudo xbps-install wofi

# Status Bar
sudo xbps-install waybar

# Notifications
sudo xbps-install mako libnotify

# Screenshot
sudo xbps-install grim slurp

# Screen Recording
sudo xbps-install wf-recorder

# Clipboard Manager
sudo xbps-install clipman
\end{lstlisting}

\section{XDG Desktop Portals (Required for Screen Sharing)}
\begin{lstlisting}
sudo xbps-install \
  xdg-desktop-portal \
  xdg-desktop-portal-wlr \
  xdg-desktop-portal-gtk \
  xdg-utils
\end{lstlisting}

\subsection{Portal Configuration File (Important!)}
\begin{lstlisting}
mkdir -p ~/.config/xdg-desktop-portal
nano ~/.config/xdg-desktop-portal/sway-portals.conf
\end{lstlisting}

\begin{lstlisting}
[preferred]
default=gtk
org.freedesktop.impl.portal.Screenshot=wlr
org.freedesktop.impl.portal.ScreenCast=wlr
\end{lstlisting}

\begin{notebox}
The file must be named \texttt{sway-portals.conf} exactly, because \texttt{XDG\_CURRENT\_DESKTOP=sway}.
\end{notebox}

\section{Media Codecs}
\begin{lstlisting}
sudo xbps-install \
  ffmpeg \
  gst-plugins-base1 \
  gst-plugins-good1 \
  gst-plugins-bad1 \
  gst-plugins-ugly1
\end{lstlisting}

\section{Power Management (Laptops)}
\begin{lstlisting}
sudo xbps-install tlp acpid acpi brightnessctl
sudo ln -s /etc/sv/acpid /var/service/
sudo ln -s /etc/sv/tlpd /var/service/
sudo usermod -aG video $USER
\end{lstlisting}

\subsection{TLP Configuration (Important for WiFi stability)}
\begin{lstlisting}
sudo nano /etc/tlp.conf
\end{lstlisting}
Find and change:
\begin{lstlisting}
WIFI_PWR_ON_AC=off
WIFI_PWR_ON_BAT=off
\end{lstlisting}

\section{Fonts}
\begin{lstlisting}
sudo xbps-install \
  liberation-fonts-ttf \
  dejavu-fonts-ttf \
  noto-fonts-ttf \
  noto-fonts-ttf-extra \
  noto-fonts-emoji \
  font-ibm-plex-ttf \
  terminus-font \
  font-awesome6
\end{lstlisting}

\begin{notebox}
\textbf{JetBrains Mono:} Not in the repo --- download it manually from \url{https://www.jetbrains.com/lp/mono/} and place it in \texttt{\textasciitilde/.local/share/fonts/}
\end{notebox}

\section{Environment Variables}

Add to \texttt{\textasciitilde/.bash\_profile} or \texttt{\textasciitilde/.profile}:

\begin{lstlisting}
# XDG Base Directories
export XDG_CONFIG_HOME="$HOME/.config"
export XDG_DATA_HOME="$HOME/.local/share"
export XDG_CACHE_HOME="$HOME/.cache"
export XDG_STATE_HOME="$HOME/.local/state"
export PATH="$HOME/.local/bin:$PATH"

# XDG_RUNTIME_DIR (if not using elogind)
if test -z "${XDG_RUNTIME_DIR}"; then
    export XDG_RUNTIME_DIR=/tmp/${UID}-runtime-dir
    if ! test -d "${XDG_RUNTIME_DIR}"; then
        mkdir "${XDG_RUNTIME_DIR}"
        chmod 0700 "${XDG_RUNTIME_DIR}"
    fi
fi

# Wayland
export XDG_SESSION_TYPE=wayland
export XDG_SESSION_DESKTOP=sway
export XDG_CURRENT_DESKTOP=sway
export DE=sway

# Application hints
export MOZ_ENABLE_WAYLAND=1
export QT_QPA_PLATFORM=wayland
export QT_WAYLAND_DISABLE_WINDOWDECORATION=1
export SDL_VIDEODRIVER=wayland
export ELECTRON_OZONE_PLATFORM_HINT=auto
export _JAVA_AWT_WM_NONREPARENTING=1

# If using seatd
export LIBSEAT_BACKEND=seatd

# Auto-start Sway from tty1
if [ -z "$DISPLAY" ] && [ -z "$WAYLAND_DISPLAY" ] && [ "${XDG_VTNR:-0}" -eq 1 ]; then
    exec dbus-run-session sway
fi
\end{lstlisting}

\section{Copy the Default Config}
\begin{lstlisting}
mkdir -p ~/.config/sway
cp /etc/sway/config ~/.config/sway/config
\end{lstlisting}

\section{Essential Additions to the Sway Config}

Open \texttt{\textasciitilde/.config/sway/config} and add:

\begin{lstlisting}
### Update dbus environment ###
exec dbus-update-activation-environment DISPLAY WAYLAND_DISPLAY SWAYSOCK XDG_CURRENT_DESKTOP

### XDG Desktop Portal ###
exec /usr/libexec/xdg-desktop-portal-wlr &
exec /usr/libexec/xdg-desktop-portal -r &

### PipeWire ###
exec pipewire &

### Notifications ###
exec mako &

### Polkit ###
exec /usr/libexec/polkit-gnome-authentication-agent-1 &

### Network ###
exec nm-applet --indicator &

### Bluetooth ###
exec blueman-applet &

### Clipboard ###
exec wl-paste -t text --watch clipman store &

### USB Auto-mount ###
exec udiskie --tray &

### Keyboard Layout ###
input type:keyboard {
    xkb_layout us,ara
    xkb_options grp:lwin_toggle
    repeat_delay 300
    repeat_rate 50
}

### Touchpad ###
input type:touchpad {
    tap enabled
    natural_scroll enabled
    dwt enabled
}

### Key Bindings ###
# Screenshot
bindsym Print exec grim ~/Pictures/screenshot-$(date +%Y%m%d-%H%M%S).png
bindsym Shift+Print exec grim -g "$(slurp)" ~/Pictures/screenshot-$(date +%Y%m%d-%H%M%S).png
bindsym Ctrl+Print exec grim -g "$(slurp)" - | wl-copy

# Volume
bindsym XF86AudioRaiseVolume exec pactl set-sink-volume @DEFAULT_SINK@ +5%
bindsym XF86AudioLowerVolume exec pactl set-sink-volume @DEFAULT_SINK@ -5%
bindsym XF86AudioMute exec pactl set-sink-mute @DEFAULT_SINK@ toggle

# Brightness
bindsym XF86MonBrightnessUp exec brightnessctl set +10%
bindsym XF86MonBrightnessDown exec brightnessctl set 10%-

# Lock screen
bindsym $mod+l exec swaylock -f -c 000000

### Idle + Lock ###
exec swayidle -w \
    timeout 300 'swaylock -f -c 000000' \
    timeout 600 'swaymsg "output * dpms off"' \
    resume 'swaymsg "output * dpms on"' \
    before-sleep 'swaylock -f -c 000000'
\end{lstlisting}

\section{Final Verification}
\begin{lstlisting}
sv status dbus
sv status elogind   # or seatd
groups $USER
echo $XDG_RUNTIME_DIR
echo $WAYLAND_DISPLAY   # after launching Sway: wayland-1
pactl info | grep "Server Name"
grim ~/test.png
\end{lstlisting}

% =============================================================
\chapter{Audio (PipeWire + ALSA)}
% =============================================================

\section{Installation}
\begin{lstlisting}
sudo xbps-install \
  pipewire \
  wireplumber \
  alsa-utils \
  alsa-tools \
  alsa-pipewire \
  libspa-bluetooth \
  pavucontrol
\end{lstlisting}

\begin{notebox}
\texttt{libspa-bluetooth} is required for Bluetooth headphones to work with PipeWire.
\end{notebox}

\section{Configure WirePlumber}
\begin{lstlisting}
mkdir -p ~/.config/pipewire/pipewire.conf.d

ln -s /usr/share/examples/wireplumber/10-wireplumber.conf \
      ~/.config/pipewire/pipewire.conf.d/
\end{lstlisting}

\section{Configure PipeWire-Pulse (for PulseAudio app support)}
\begin{lstlisting}
sudo xbps-install pipewire-pulse

ln -s /usr/share/examples/pipewire/20-pipewire-pulse.conf \
      ~/.config/pipewire/pipewire.conf.d/
\end{lstlisting}

\section{Configure ALSA with PipeWire}
\begin{lstlisting}
sudo mkdir -p /etc/alsa/conf.d

sudo ln -s /usr/share/alsa/alsa.conf.d/50-pipewire.conf \
           /etc/alsa/conf.d/

sudo ln -s /usr/share/alsa/alsa.conf.d/99-pipewire-default.conf \
           /etc/alsa/conf.d/

alsactl init
\end{lstlisting}

\section{Starting PipeWire}

On Void, PipeWire is \textbf{started from inside Sway}, not as a system service.

In your Sway config:
\begin{lstlisting}
exec pipewire &
\end{lstlisting}

\begin{warningbox}
\textbf{Warning:} Do not try to enable it with \texttt{ln -s /etc/sv/pipewire /var/service/} --- it will not work correctly.
\end{warningbox}

\section{Adjusting Volume Levels}
\begin{lstlisting}
# Command line
alsamixer

# Graphical
pavucontrol
\end{lstlisting}

\section{Verification}
\begin{lstlisting}
pactl info | grep "Server Name"
# Expected: PulseAudio (on PipeWire ...)

pactl list sources short
pactl list sinks short
\end{lstlisting}

\section{Common Problems}

\textbf{No audio:}
\begin{lstlisting}
pactl info
pipewire &
pactl info
\end{lstlisting}

\textbf{Bluetooth headphones not working:}
\begin{lstlisting}
sudo xbps-install libspa-bluetooth
# Then restart Sway
\end{lstlisting}

\textbf{No sound device found:}
\begin{lstlisting}
aplay -l
# If empty, check GPU drivers and kernel modules
\end{lstlisting}

% =============================================================
\chapter{Theming (Consistent Look)}
% =============================================================

\section{The Basic Concept}

A Wayland theme consists of 4 separate layers:

\begin{center}
\begin{tabularx}{0.85\textwidth}{lL}
\toprule
\textbf{Layer} & \textbf{Effect} \\
\midrule
GTK Theme   & Colors and shape of buttons and windows \\
Icon Theme  & File and application icons \\
Cursor Theme & Mouse pointer appearance \\
Font        & Text font in all applications \\
\bottomrule
\end{tabularx}
\end{center}

\begin{notebox}
\textbf{The problem:} On Wayland, GTK does not read settings automatically --- they must be set via \texttt{gsettings}. \textbf{The solution:} Use the \texttt{nwg-look} tool.
\end{notebox}

\section{Installation}
\begin{lstlisting}
sudo xbps-install nwg-look

# Materia (flat, modern)
sudo xbps-install materia-theme

# Arc
sudo xbps-install arc-theme

# Adwaita:dark -- included with GTK by default
\end{lstlisting}

Download a theme from the internet (e.g.\ Catppuccin):
\begin{lstlisting}
mkdir -p ~/.local/share/themes
cd ~/.local/share/themes
# Extract the downloaded zip here
\end{lstlisting}

\begin{lstlisting}
# Icon Theme
sudo xbps-install papirus-icon-theme

# Cursor Theme
sudo xbps-install breeze-cursors
\end{lstlisting}

\section{Applying the Theme}
\begin{lstlisting}
nwg-look
# Choose GTK Theme + Icons + Cursor + Font
# Click Apply
\end{lstlisting}

\section{Cursor in the Sway Config}

Add to \texttt{\textasciitilde/.config/sway/config}:
\begin{lstlisting}
seat seat0 xcursor_theme THEME_NAME 24
\end{lstlisting}

\section{Fonts}
\begin{lstlisting}
sudo xbps-install \
  font-jetbrains-mono \
  noto-fonts-ttf \
  noto-fonts-ttf-extra \
  noto-fonts-emoji \
  font-awesome6 \
  dejavu-fonts-ttf

# Rebuild cache after installing new fonts
fc-cache -fv
\end{lstlisting}

In the Sway config:
\begin{lstlisting}
font pango:JetBrains Mono 11, Font Awesome 6 Free 11
\end{lstlisting}

\section{Theme Locations}
\begin{lstlisting}
# GTK themes
/usr/share/themes/
~/.local/share/themes/

# Icons + Cursors
/usr/share/icons/
~/.local/share/icons/

# List installed themes
ls /usr/share/themes/
ls /usr/share/icons/
fc-list | grep -i "jet\|noto"
\end{lstlisting}

% =============================================================
\chapter{Screenshot, Screen Recording, and Clipboard}
% =============================================================

\section{Install the Tools}
\begin{lstlisting}
sudo xbps-install grim slurp wf-recorder wl-clipboard clipman
\end{lstlisting}

\section{Screenshots}

\subsection{Direct Commands}
\begin{lstlisting}
# Full screen
grim ~/Pictures/screenshot-$(date +%Y%m%d-%H%M%S).png

# Selected region
grim -g "$(slurp)" ~/Pictures/screenshot-$(date +%Y%m%d-%H%M%S).png

# Copy directly to clipboard
grim -g "$(slurp)" - | wl-copy
\end{lstlisting}

\subsection{In the Sway Config}
\begin{lstlisting}
bindsym Print exec grim ~/Pictures/screenshot-$(date +%Y%m%d-%H%M%S).png
bindsym Shift+Print exec grim -g "$(slurp)" ~/Pictures/screenshot-$(date +%Y%m%d-%H%M%S).png
bindsym Ctrl+Print exec grim -g "$(slurp)" - | wl-copy
\end{lstlisting}

\section{Screen Recording with wf-recorder}
\begin{lstlisting}
# Full screen
wf-recorder -f ~/Videos/recording.mp4

# Selected region
wf-recorder -g "$(slurp)" -f ~/Videos/recording.mp4

# With audio
wf-recorder -a -f ~/Videos/recording.mp4

# With a specific audio device
pactl list sources short          # list available devices
wf-recorder -a "source_name" -f ~/Videos/recording.mp4

# Stop recording
# Ctrl+C
\end{lstlisting}

\section{Clipboard Manager with clipman}

\subsection{Start (in Sway config)}
\begin{lstlisting}
exec wl-paste -t text --watch clipman store &
\end{lstlisting}

\subsection{Browse History}
\begin{lstlisting}
clipman pick --tool wofi
# or
clipman pick --tool fuzzel
\end{lstlisting}

Add to Sway config:
\begin{lstlisting}
bindsym $mod+shift+v exec clipman pick --tool fuzzel
\end{lstlisting}

\subsection{Clear History}
\begin{lstlisting}
clipman clear --all
\end{lstlisting}

% =============================================================
\chapter{Hardware (Bluetooth, Printer, USB, MTP)}
% =============================================================

\section{Bluetooth}

\subsection{Installation}
\begin{lstlisting}
sudo xbps-install bluez blueman libspa-bluetooth
sudo ln -s /etc/sv/bluetoothd /var/service/
sudo usermod -aG bluetooth $USER
# Log out and back in
\end{lstlisting}

In the Sway config:
\begin{lstlisting}
exec blueman-applet &
\end{lstlisting}

\subsection{Usage (CLI)}
\begin{lstlisting}
bluetoothctl
power on
scan on
pair XX:XX:XX:XX:XX:XX
connect XX:XX:XX:XX:XX:XX
trust XX:XX:XX:XX:XX:XX
\end{lstlisting}

\subsection{Verification}
\begin{lstlisting}
sv status bluetoothd
\end{lstlisting}

\section{Printer with CUPS}

\subsection{Installation}
\begin{lstlisting}
sudo xbps-install cups cups-filters gutenprint system-config-printer
# HP printers:
sudo xbps-install hplip
# Epson printers:
sudo xbps-install epson-inkjet-printer-escpr

sudo ln -s /etc/sv/cupsd /var/service/
sudo usermod -aG lpadmin $USER
\end{lstlisting}

\subsection{Adding a Printer}
Open in your browser: \texttt{http://localhost:631}\\
Then: Administration $\to$ Add Printer

\section{USB Auto-mount with udiskie}

\subsection{Installation}
\begin{lstlisting}
sudo xbps-install udisks2 udiskie gvfs ntfs-3g exfat-utils
\end{lstlisting}

\subsection{Polkit Rule (to avoid password prompts)}
\begin{lstlisting}
sudo mkdir -p /etc/polkit-1/rules.d/
sudo nano /etc/polkit-1/rules.d/50-udisks.rules
\end{lstlisting}

\begin{lstlisting}
polkit.addRule(function(action, subject) {
  if (subject.isInGroup("wheel")) {
    if (action.id.startsWith("org.freedesktop.udisks2.")) {
      return polkit.Result.YES;
    }
  }
});
\end{lstlisting}

In the Sway config:
\begin{lstlisting}
exec udiskie --tray &
\end{lstlisting}

\subsection{Manual Unmount}
\begin{lstlisting}
udiskie-umount /run/media/$USER/DRIVE_NAME
udiskie-umount --all
\end{lstlisting}

\section{MTP (Phone via USB)}
\begin{lstlisting}
sudo xbps-install go-mtpfs libmtp
sudo usermod -aG plugdev $USER
\end{lstlisting}

\subsection{FUSE Setup}
\begin{lstlisting}
# Uncomment the user_allow_other line
sudo sed -i '/user_allow_other/s/^#//g' /etc/fuse.conf
\end{lstlisting}

\subsection{Usage}
\begin{lstlisting}
# Detect connected devices
mtp-detect

# Mount the device
mkdir -p ~/phone
go-mtpfs ~/phone &

# Unmount
fusermount -u ~/phone
\end{lstlisting}

% =============================================================
\chapter{Networking (VPN, Hotspot, Tethering)}
% =============================================================

\section{NetworkManager Setup}
\begin{lstlisting}
sudo xbps-install NetworkManager network-manager-applet

# Stop old services
sudo sv stop dhcpcd wpa_supplicant 2>/dev/null
sudo rm /var/service/dhcpcd /var/service/wpa_supplicant 2>/dev/null

# Enable NetworkManager
sudo ln -s /etc/sv/dbus /var/service/
sudo ln -s /etc/sv/NetworkManager /var/service/
\end{lstlisting}

\subsection{Static DNS (optional)}
\begin{lstlisting}
sudo nano /etc/resolv.conf
\end{lstlisting}
\begin{lstlisting}
nameserver 1.1.1.1
nameserver 8.8.8.8
\end{lstlisting}
\begin{lstlisting}
sudo chattr +i /etc/resolv.conf   # Prevent NetworkManager from overwriting it
\end{lstlisting}

\section{VPN}

\subsection{Install Plugins}
\begin{lstlisting}
# OpenVPN
sudo xbps-install NetworkManager-openvpn

# WireGuard
sudo xbps-install wireguard-tools NetworkManager-wireguard

# L2TP/IPSec
sudo xbps-install NetworkManager-l2tp
\end{lstlisting}

\subsection{Add a VPN Connection}
\begin{lstlisting}
# Import a config file
nmcli connection import type openvpn file /path/to/config.ovpn
nmcli connection import type wireguard file /path/to/config.conf

# Using the TUI
nmtui
\end{lstlisting}

\subsection{Enable / Disable}
\begin{lstlisting}
nmcli connection up VPN_NAME
nmcli connection down VPN_NAME
nmcli connection show
\end{lstlisting}

\section{Hotspot}

\begin{warningbox}
\textbf{Warning:} Most WiFi cards do not support internet access and hotspot simultaneously on the same adapter.
\end{warningbox}

\subsection{Quick Start}
\begin{lstlisting}
nmcli device wifi hotspot ifname wlan0 ssid "MyNetwork" password "MyPassword"
\end{lstlisting}

\subsection{Show QR Code}
\begin{lstlisting}
nmcli device wifi show-password
\end{lstlisting}

\subsection{Persistent Hotspot}
\begin{lstlisting}
nmcli con add type wifi ifname wlan0 con-name "MyHotspot" autoconnect no \
  ssid "MyNetwork" \
  802-11-wireless.mode ap \
  802-11-wireless.band bg \
  ipv4.method shared \
  wifi-sec.key-mgmt wpa-psk \
  wifi-sec.psk "MyPassword"

nmcli con up MyHotspot
nmcli con down MyHotspot
\end{lstlisting}

\section{USB Tethering (Phone Internet on Laptop)}
\begin{lstlisting}
# Connect your phone via USB and enable USB Tethering on it
ip link show             # A new interface will appear (usb0 or enp...)
nmcli device status      # NetworkManager should detect it automatically
\end{lstlisting}

% =============================================================
\chapter{Applications}
% =============================================================

\section{Browser}
\begin{lstlisting}
sudo xbps-install firefox
\end{lstlisting}
Firefox runs natively on Wayland as long as \texttt{MOZ\_ENABLE\_WAYLAND=1} is set in your environment.

\section{Multimedia}
\begin{lstlisting}
# Video
sudo xbps-install mpv        # native Wayland
# or
sudo xbps-install vlc

# Images
sudo xbps-install imv         # native Wayland

# Audio control
sudo xbps-install pavucontrol
\end{lstlisting}

\section{File Manager}
\begin{lstlisting}
# Terminal
sudo xbps-install nnn

# Graphical
sudo xbps-install thunar
\end{lstlisting}

\section{Archive Tools}
\begin{lstlisting}
sudo xbps-install unzip zip p7zip file-roller
\end{lstlisting}

\section{System Monitoring}
\begin{lstlisting}
sudo xbps-install htop btop fastfetch lm_sensors
\end{lstlisting}

\section{Flatpak}
\begin{lstlisting}
sudo xbps-install flatpak xdg-user-dirs
flatpak remote-add --if-not-exists flathub https://dl.flathub.org/repo/flathub.flatpakrepo

# Install an app
flatpak install flathub com.discordapp.Discord

# Run an app
flatpak run com.discordapp.Discord

# Update all
flatpak update

# List installed apps
flatpak list
\end{lstlisting}

\section{AppImage}
\begin{lstlisting}
sudo xbps-install fuse
sudo modprobe fuse

chmod +x AppName.AppImage
./AppName.AppImage

# If it complains about FUSE
./AppName.AppImage --appimage-extract-and-run

# Organized location
mkdir -p ~/AppImages
\end{lstlisting}

\section{Emacs}
\begin{lstlisting}
sudo xbps-install emacs        # with GUI
sudo xbps-install emacs-nox    # terminal only (lighter)
\end{lstlisting}
Config location: \texttt{\textasciitilde/.config/emacs/init.el}

\section{Terminal Tools for Coding}
\begin{lstlisting}
sudo xbps-install \
  ripgrep \
  fd \
  bat \
  eza \
  fzf \
  zoxide \
  btop \
  lazygit
\end{lstlisting}

\begin{center}
\begin{tabularx}{0.9\textwidth}{llL}
\toprule
\textbf{Tool} & \textbf{Replaces} & \textbf{Benefit} \\
\midrule
rg      & grep  & Much faster \\
fd      & find  & Simpler and faster \\
bat     & cat   & Syntax highlighting \\
eza     & ls    & Cleaner and more informative \\
fzf     & ---   & Fuzzy finder for everything \\
zoxide  & cd    & Remembers visited directories \\
btop    & htop  & More features and better visuals \\
lazygit & git   & TUI for git \\
\bottomrule
\end{tabularx}
\end{center}

\subsection{Link fzf to Shell History}
Add to \texttt{\textasciitilde/.bashrc}:
\begin{lstlisting}
source /usr/share/fzf/key-bindings.bash
source /usr/share/fzf/completion.bash
\end{lstlisting}
Then \texttt{Ctrl+R} will open a fuzzy search over your shell history.

% =============================================================
\chapter{AppImage and Desktop Entries}
% =============================================================

\section{Directory Structure}
\begin{lstlisting}
~/AppImages/
 +-- obsidian.AppImage
 +-- ticktick.AppImage
 +-- icons/
     +-- obsidian.png
     +-- ticktick.png

~/.local/share/applications/    # Desktop entries
~/.config/mimeapps.list         # MIME associations
\end{lstlisting}

\section{Setting Up a New AppImage}
\begin{lstlisting}
# 1. Create directories
mkdir -p ~/AppImages/icons
mkdir -p ~/.local/share/applications

# 2. Move the AppImage
mv downloaded.AppImage ~/AppImages/appname.AppImage
chmod 755 ~/AppImages/appname.AppImage

# 3. Place the icon
cp icon.png ~/AppImages/icons/appname.png
chmod 644 ~/AppImages/icons/appname.png
\end{lstlisting}

\section{Desktop Entry Template}
\begin{lstlisting}
nano ~/.local/share/applications/appname.desktop
\end{lstlisting}

\begin{lstlisting}
[Desktop Entry]
Version=1.0
Type=Application
Name=Application Name
Comment=Short description
Exec=/home/USERNAME/AppImages/appname.AppImage %F
Icon=/home/USERNAME/AppImages/icons/appname.png
Terminal=false
Categories=Office;Utility;
MimeType=text/markdown;text/plain;
StartupNotify=true
\end{lstlisting}

\begin{lstlisting}
chmod 644 ~/.local/share/applications/appname.desktop
update-desktop-database ~/.local/share/applications
\end{lstlisting}

\section{Exec Variables}

\begin{center}
\begin{tabularx}{0.7\textwidth}{lL}
\toprule
\textbf{Variable} & \textbf{Meaning} \\
\midrule
\%f & Single file \\
\%F & Multiple files \\
\%u & Single URL \\
\%U & Multiple URLs \\
\bottomrule
\end{tabularx}
\end{center}

\section{Common MIME Types}
\begin{lstlisting}
text/plain        ->  .txt
text/markdown     ->  .md
text/html         ->  .html
application/pdf   ->  .pdf
image/png         ->  .png
image/jpeg        ->  .jpg
\end{lstlisting}

\subsection{Useful Commands}
\begin{lstlisting}
# Get the MIME type of a file
xdg-mime query filetype document.md

# Get the default application for a type
xdg-mime query default text/markdown

# Set the default application
xdg-mime default obsidian.desktop text/markdown
\end{lstlisting}

\section{Full Example: Obsidian}
\begin{lstlisting}
mkdir -p ~/AppImages/icons
mv Obsidian-*.AppImage ~/AppImages/obsidian.AppImage
chmod 755 ~/AppImages/obsidian.AppImage
cp obsidian-icon.png ~/AppImages/icons/obsidian.png
chmod 644 ~/AppImages/icons/obsidian.png
\end{lstlisting}

\begin{lstlisting}
# ~/.local/share/applications/obsidian.desktop
[Desktop Entry]
Version=1.0
Type=Application
Name=Obsidian
Comment=Knowledge base and note-taking app
Exec=/home/USERNAME/AppImages/obsidian.AppImage %F
Icon=/home/USERNAME/AppImages/icons/obsidian.png
Terminal=false
Categories=Office;TextEditor;Utility;
MimeType=text/markdown;text/plain;
StartupNotify=true
\end{lstlisting}

\begin{lstlisting}
chmod 644 ~/.local/share/applications/obsidian.desktop
update-desktop-database ~/.local/share/applications
xdg-mime default obsidian.desktop text/markdown
\end{lstlisting}

\section{Web App as a Desktop Entry}
\begin{lstlisting}
[Desktop Entry]
Version=1.0
Type=Application
Name=Gmail
Exec=firefox --new-window https://mail.google.com
Icon=/home/USERNAME/AppImages/icons/gmail.png
Terminal=false
Categories=Network;WebBrowser;
StartupNotify=true
\end{lstlisting}

\section{Quick Reference}
\begin{lstlisting}
chmod +x file.AppImage
chmod 644 file.desktop
update-desktop-database ~/.local/share/applications
xdg-mime default app.desktop text/markdown
xdg-mime query default text/markdown
gtk-launch appname              # Test a desktop entry
xdg-open file.md                # Open with the default application
desktop-file-validate app.desktop
\end{lstlisting}

% =============================================================
\chapter{KVM/QEMU Virtualization}
% =============================================================

\section{Verify CPU Virtualization Support}
\begin{lstlisting}
lscpu | grep -i virtualization
grep -E '(vmx|svm)' /proc/cpuinfo
\end{lstlisting}

\begin{itemize}[leftmargin=*]
  \item \textbf{Intel:} Must show \texttt{VT-x} or \texttt{vmx}
  \item \textbf{AMD:} Must show \texttt{AMD-V} or \texttt{svm}
\end{itemize}

If nothing appears, enable it in the BIOS: look for \textit{Virtualization Technology} or \textit{Intel VT-x}.

\section{Installation}
\begin{lstlisting}
sudo xbps-install virt-manager libvirt qemu dnsmasq bridge-utils openbsd-netcat
\end{lstlisting}

\begin{center}
\begin{tabularx}{0.85\textwidth}{lL}
\toprule
\textbf{Package} & \textbf{Purpose} \\
\midrule
virt-manager  & GUI for managing VMs \\
libvirt       & Virtualization API \\
qemu          & The emulator \\
dnsmasq       & DNS/DHCP for virtual networks \\
bridge-utils  & Network bridging \\
\bottomrule
\end{tabularx}
\end{center}

\section{Enable Services}
\begin{lstlisting}
sudo ln -s /etc/sv/dbus /var/service/        # if not already enabled
sudo ln -s /etc/sv/libvirtd /var/service/
sudo ln -s /etc/sv/virtlockd /var/service/
sudo ln -s /etc/sv/virtlogd /var/service/
\end{lstlisting}

\section{Add User to the Group}
\begin{lstlisting}
sudo usermod -aG libvirt $USER
# Log out and back in
\end{lstlisting}

\section{Load the KVM Module}
\begin{lstlisting}
# Intel
sudo modprobe kvm_intel

# AMD
sudo modprobe kvm_amd

# Load automatically on boot (Intel)
echo "kvm_intel" | sudo tee /etc/modules-load.d/kvm.conf
\end{lstlisting}

\section{Configure libvirt (to run VMs without sudo)}
\begin{lstlisting}
sudo nano /etc/libvirt/libvirtd.conf
\end{lstlisting}

Uncomment these lines:
\begin{lstlisting}
unix_sock_group = "libvirt"
unix_sock_ro_perms = "0777"
unix_sock_rw_perms = "0770"
\end{lstlisting}

\begin{lstlisting}
sudo nano /etc/libvirt/libvirt.conf
\end{lstlisting}

Uncomment:
\begin{lstlisting}
uri_default = "qemu:///system"
\end{lstlisting}

\begin{lstlisting}
sudo sv restart libvirtd
sudo reboot
\end{lstlisting}

\section{Post-Reboot Verification}
\begin{lstlisting}
lsmod | grep kvm
ls -la /dev/kvm
sudo virsh list --all
groups | grep libvirt
\end{lstlisting}

\section{Essential virsh Commands}
\begin{lstlisting}
virsh list --all              # List all VMs
virsh start <vm-name>         # Start a VM
virsh shutdown <vm-name>      # Graceful shutdown
virsh destroy <vm-name>       # Force shutdown
virsh undefine <vm-name>      # Delete a VM
virsh dominfo <vm-name>       # Show VM info
virsh edit <vm-name>          # Edit VM config
\end{lstlisting}

\section{Nested Virtualization}
\begin{lstlisting}
# Intel
sudo modprobe -r kvm_intel
sudo modprobe kvm_intel nested=1

# Persistent
echo "options kvm_intel nested=1" | sudo tee /etc/modprobe.d/kvm.conf
\end{lstlisting}

\section{Troubleshooting}
\begin{lstlisting}
# "Unable to connect to libvirt"
sudo sv status dbus
sudo sv status libvirtd
sudo sv up libvirtd

# "Could not access KVM kernel module"
lsmod | grep kvm
sudo modprobe kvm_intel
ls -la /dev/kvm

# "Permission denied"
groups | grep libvirt
sudo usermod -aG libvirt $USER
# Log out and back in
\end{lstlisting}

% =============================================================
\chapter{Troubleshooting}
% =============================================================

\section{Sway Won't Start}

\subsection{"Failed to create backend"}
\begin{lstlisting}
# Verify elogind or seatd
sudo sv status elogind
sudo sv status seatd

# If using seatd, verify group membership
groups | grep _seatd

# Verify dbus
sudo sv status dbus
\end{lstlisting}

\subsection{Sway starts and crashes immediately}
\begin{lstlisting}
dbus-run-session sway 2>&1 | head -30
\end{lstlisting}

\section{XDG\_RUNTIME\_DIR Issues}
\begin{lstlisting}
# Temporary manual fix
export XDG_RUNTIME_DIR=/run/user/$(id -u)
sudo mkdir -p $XDG_RUNTIME_DIR
sudo chown $USER:$USER $XDG_RUNTIME_DIR
sudo chmod 0700 $XDG_RUNTIME_DIR

# Permanent fix: pam_rundir (with seatd)
sudo xbps-install pam_rundir
# Add to /etc/pam.d/system-login:
# -session optional pam_rundir.so
\end{lstlisting}

\section{Audio}

\subsection{PipeWire not running}
\begin{lstlisting}
sudo sv status dbus
ls ~/.config/pipewire/pipewire.conf.d/
pipewire &
pactl info
grep pipewire ~/.config/sway/config
\end{lstlisting}

\section{Firefox Not Running on Wayland}
\begin{lstlisting}
echo $MOZ_ENABLE_WAYLAND
# Expected: 1

MOZ_ENABLE_WAYLAND=1 firefox
\end{lstlisting}

\section{Screen Sharing Not Working}
\begin{lstlisting}
# 1. Verify installation
xbps-query xdg-desktop-portal-wlr

# 2. Check portal config
cat ~/.config/xdg-desktop-portal/sway-portals.conf

# 3. Check dbus environment variables
dbus-update-activation-environment --print-environment | grep XDG_CURRENT_DESKTOP

# 4. Restart portals manually
killall xdg-desktop-portal xdg-desktop-portal-wlr 2>/dev/null
/usr/libexec/xdg-desktop-portal-wlr &
/usr/libexec/xdg-desktop-portal -r &

# 5. Test in Firefox
# https://mozilla.github.io/webrtc-landing/gum_test.html
\end{lstlisting}

\section{Qt Apps Showing as X11}
\begin{lstlisting}
echo $QT_QPA_PLATFORM
# Expected: wayland

sudo xbps-install qt5-wayland qt6-wayland
\end{lstlisting}

\section{Java Apps (Blank Grey Window)}
\begin{lstlisting}
# Add to ~/.profile
export _JAVA_AWT_WM_NONREPARENTING=1
\end{lstlisting}

\section{AppImage Won't Launch}
\begin{lstlisting}
sudo xbps-install fuse
sudo modprobe fuse
chmod 755 ~/AppImages/app.AppImage
./app.AppImage --appimage-extract-and-run
\end{lstlisting}

\section{General Diagnostic Commands}
\begin{lstlisting}
echo $WAYLAND_DISPLAY        # Expected: wayland-1
swaymsg -t get_version
swaymsg -t get_outputs
swaymsg -t get_inputs
dbus-update-activation-environment --print-environment
\end{lstlisting}

% =============================================================
\chapter{xbps-src: Building Packages Not in the Repo}
% =============================================================

\section{The Core Concept}

\texttt{xbps-src} is similar to \texttt{PKGBUILD} on Arch Linux or \texttt{ebuild} on Gentoo.

Instead of building from source manually, you write (or copy) a small ``recipe'' file called a \textbf{template}. Then \texttt{xbps-src} builds the program and creates an \texttt{.xbps} package that you can install and remove with \texttt{xbps} like any other package.

\begin{notebox}
\textbf{The main benefit:} The program becomes managed --- you can remove it with \texttt{xbps-remove} just like any other package.
\end{notebox}

\section{Initial Setup (One Time Only)}

\begin{lstlisting}
sudo xbps-install git xtools base-devel

cd ~
git clone --depth=1 https://github.com/void-linux/void-packages.git
cd void-packages

./xbps-src binary-bootstrap
\end{lstlisting}

\begin{notebox}
\texttt{--depth=1} fetches only the latest snapshot without the full history --- much faster.
\end{notebox}

\section{Building a Package Already in void-packages}

\begin{lstlisting}
# Search for the template
ls srcpkgs/ | grep -i package_name
./xbps-src search package_name

# Build
./xbps-src pkg package_name

# Install
sudo xbps-install --repository=hostdir/binpkgs package_name

# If in nonfree
sudo xbps-install --repository=hostdir/binpkgs/nonfree package_name
\end{lstlisting}

\subsection{Practical Example: Google Chrome}
\begin{lstlisting}
cd ~/void-packages
echo "XBPS_ALLOW_RESTRICTED=yes" >> etc/conf
./xbps-src pkg google-chrome
sudo xbps-install --repository=hostdir/binpkgs/nonfree google-chrome
\end{lstlisting}

\section{Adding the Local Repo Permanently}
\begin{lstlisting}
sudo nano /etc/xbps.d/local-repo.conf
\end{lstlisting}

\begin{lstlisting}
repository=/home/YOUR_USERNAME/void-packages/hostdir/binpkgs
repository=/home/YOUR_USERNAME/void-packages/hostdir/binpkgs/nonfree
\end{lstlisting}

\section{Anatomy of a Template}

\begin{lstlisting}
pkgname=my-tool
version=1.2.3
revision=1
short_desc="An awesome tool"
maintainer="Your Name <email@example.com>"
license="MIT"
homepage="https://github.com/someone/my-tool"

distfiles="https://github.com/someone/my-tool/archive/v${version}.tar.gz"
checksum="sha256 hash goes here"

hostmakedepends="python3"
depends="python3"

do_install() {
    python3 setup.py install --prefix=/usr --root="${DESTDIR}"
}
\end{lstlisting}

\subsection{Key Variables}

\begin{center}
\begin{tabularx}{0.9\textwidth}{lL}
\toprule
\textbf{Variable} & \textbf{Meaning} \\
\midrule
pkgname        & Package name \\
version        & Version number \\
revision       & Revision number (start at 1) \\
short\_desc    & Short description (max 72 characters) \\
distfiles      & Source URL (tar.gz or zip) \\
checksum       & SHA256 hash of the source file \\
hostmakedepends & Build-time tools \\
makedepends    & Build-time libraries \\
depends        & Runtime dependencies \\
build\_style   & Automatic build method \\
\bottomrule
\end{tabularx}
\end{center}

\subsection{build\_style Options}
\begin{lstlisting}
build_style=cmake          # CMake projects
build_style=meson          # Meson projects
build_style=gnu-configure  # Autotools (./configure)
build_style=python3-module # Python (setup.py)
build_style=go             # Go
build_style=cargo          # Rust
build_style=gnu-makefile   # Simple Makefile projects
\end{lstlisting}

\section{Getting the Checksum}
\begin{lstlisting}
# Using xtools
xgensum -i srcpkgs/package_name/template

# Manually
wget https://source_url
sha256sum filename.tar.gz
\end{lstlisting}

\section{Full Example: A Rust Package (zoxide)}
\begin{lstlisting}
./xbps-src newpkg zoxide
\end{lstlisting}

\begin{lstlisting}
# srcpkgs/zoxide/template
pkgname=zoxide
version=0.9.4
revision=1
build_style=cargo
short_desc="Smarter cd command"
maintainer="Me <me@example.com>"
license="MIT"
homepage="https://github.com/ajeetdsouza/zoxide"
distfiles="https://github.com/ajeetdsouza/zoxide/archive/v${version}.tar.gz"
checksum="put sha256 here"
\end{lstlisting}

\begin{lstlisting}
xgensum -i srcpkgs/zoxide/template
./xbps-src pkg zoxide
sudo xbps-install --repository=hostdir/binpkgs zoxide
\end{lstlisting}

\section{Everyday xbps-src Commands}
\begin{lstlisting}
./xbps-src pkg package_name          # Build
./xbps-src -f pkg package_name       # Force rebuild from scratch
./xbps-src clean package_name        # Clean build files
./xbps-src show package_name         # Show dependencies
cd ~/void-packages && git pull        # Update void-packages
./xbps-src update-sys                 # Update outdated packages
\end{lstlisting}

\section{Updating a Package}
\begin{lstlisting}
cd ~/void-packages
git pull
nano srcpkgs/package_name/template   # Change version=
xgensum -i srcpkgs/package_name/template
./xbps-src pkg package_name
sudo xbps-install -u package_name
\end{lstlisting}

\section{Hold: Prevent a Package from Updating}
\begin{lstlisting}
sudo xbps-pkgdb -m hold package_name
sudo xbps-pkgdb -m unhold package_name
\end{lstlisting}

\section{Troubleshooting xbps-src}

\textbf{"checksum mismatch":}
\begin{lstlisting}
xgensum -i srcpkgs/package_name/template
\end{lstlisting}

\textbf{Build is very slow:}
\begin{lstlisting}
echo "XBPS_MAKEJOBS=$(nproc)" >> etc/conf
\end{lstlisting}

\textbf{View the full build log:}
\begin{lstlisting}
./xbps-src pkg package_name 2>&1 | tee /tmp/build.log
less /tmp/build.log
\end{lstlisting}

\section{Quick Summary}
\begin{lstlisting}
# Initial setup (one time)
git clone --depth=1 https://github.com/void-linux/void-packages.git
cd void-packages
./xbps-src binary-bootstrap

# Build and install an existing package
./xbps-src pkg package_name
sudo xbps-install --repository=hostdir/binpkgs package_name

# Create a new package
./xbps-src newpkg package_name
# Edit the template
xgensum -i srcpkgs/package_name/template
./xbps-src pkg package_name
sudo xbps-install --repository=hostdir/binpkgs package_name

# Remove
sudo xbps-remove package_name
\end{lstlisting}

% =============================================================
\chapter{i3 on Void Linux (Migration from Sway + Running Both)}
% =============================================================

\section{The Core Difference in One Sentence}

\begin{center}
\begin{tabularx}{0.85\textwidth}{lLL}
\toprule
 & \textbf{Sway} & \textbf{i3} \\
\midrule
Display Server & Wayland & X11 / Xorg \\
Config syntax  & \textasciitilde99\% identical & \textasciitilde99\% identical \\
Compositor     & Built-in & Requires separate \texttt{picom} \\
Screen Sharing & Via Portal + PipeWire & Direct (easier) \\
\bottomrule
\end{tabularx}
\end{center}

\section{Full Migration Table: Swap These for Those}

\begin{center}
\begin{tabularx}{\textwidth}{lLL}
\toprule
\textbf{Role} & \textbf{Sway (Wayland)} & \textbf{i3 (X11)} \\
\midrule
Window Manager    & sway                         & i3 \\
Status Bar        & waybar                       & polybar or i3status \\
Notifications     & mako                         & dunst \\
Launcher          & wofi / fuzzel                & rofi / dmenu \\
Lock Screen       & swaylock                     & i3lock \\
Idle              & swayidle                     & xss-lock + xautolock \\
Wallpaper         & swaybg                       & feh / nitrogen \\
Screenshot        & grim + slurp                 & flameshot / scrot \\
Screen Recording  & wf-recorder                  & OBS / ffmpeg \\
Clipboard         & wl-clipboard + clipman       & xclip / xdotool \\
Compositor        & Built into Sway              & picom \\
Display Config    & kanshi / wdisplays           & arandr / autorandr \\
Brightness        & \multicolumn{2}{l}{brightnessctl (same)} \\
Volume            & \multicolumn{2}{l}{pactl (same)} \\
Theming           & nwg-look                     & lxappearance \\
Polkit Agent      & \multicolumn{2}{l}{polkit-gnome (same)} \\
Keyring           & \multicolumn{2}{l}{gnome-keyring (same)} \\
Network Applet    & \multicolumn{2}{l}{nm-applet (same)} \\
Bluetooth Applet  & \multicolumn{2}{l}{blueman-applet (same)} \\
\bottomrule
\end{tabularx}
\end{center}

\begin{notebox}
Items marked \textbf{(same)} do not need to be changed at all.
\end{notebox}

\section{Installing i3 on Void}

\begin{lstlisting}
# i3 and X11
sudo xbps-install i3 xorg-minimal xinit

# Tool replacements
sudo xbps-install \
  polybar \
  dunst \
  i3lock \
  feh \
  picom \
  flameshot \
  xclip \
  arandr \
  lxappearance \
  xss-lock \
  xautolock
\end{lstlisting}

\section{Launching i3 from the TTY}

Create \texttt{\textasciitilde/.xinitrc}:

\begin{lstlisting}
#!/bin/sh

# dbus must be first -- without it, PipeWire and Keyring won't work
eval $(dbus-launch --sh-syntax)
export DBUS_SESSION_BUS_ADDRESS

# X11 variables
export XDG_SESSION_TYPE=x11
export XDG_SESSION_DESKTOP=i3
export XDG_CURRENT_DESKTOP=i3

# Keyring (after dbus)
eval $(gnome-keyring-daemon --start --components=gpg,pkcs11,secrets,ssh \
    | sed 's/^\(.*\)/export \1/g')
export SSH_AUTH_SOCK

# PipeWire (after dbus)
pipewire &

exec i3
\end{lstlisting}

\begin{lstlisting}
chmod +x ~/.xinitrc
\end{lstlisting}

To launch:
\begin{lstlisting}
startx
\end{lstlisting}

\begin{warningbox}
\textbf{Common mistake:} If you start PipeWire before dbus, or launch it from the i3 config instead of \texttt{.xinitrc}, audio will not work. The correct order is: \texttt{dbus} $\to$ \texttt{keyring} $\to$ \texttt{pipewire} $\to$ \texttt{i3}
\end{warningbox}

\section{Auto-start from TTY1 (Sway) or TTY2 (i3)}

In \texttt{\textasciitilde/.profile}:

\begin{lstlisting}
# TTY1 -> Sway
if [ -z "$WAYLAND_DISPLAY" ] && [ "$(tty)" = "/dev/tty1" ]; then
    exec ~/.local/bin/start-sway
fi

# TTY2 -> i3
if [ -z "$DISPLAY" ] && [ "$(tty)" = "/dev/tty2" ]; then
    exec startx
fi
\end{lstlisting}

Now: \texttt{tty1} = Sway / \texttt{tty2} = i3. Switch between them with \texttt{Ctrl+Alt+F1} and \texttt{Ctrl+Alt+F2}.

\section{i3 Config (A Full Starting Point)}

\begin{lstlisting}
mkdir -p ~/.config/i3
cp /etc/i3/config ~/.config/i3/config 2>/dev/null || \
  cp /usr/share/doc/i3/config ~/.config/i3/config 2>/dev/null || \
  i3-config-wizard
\end{lstlisting}

Add the following to \texttt{\textasciitilde/.config/i3/config}:

\begin{lstlisting}
### Startup ###
exec --no-startup-id dbus-update-activation-environment --all

exec --no-startup-id pipewire &
exec --no-startup-id dunst &
exec --no-startup-id picom -b &
exec --no-startup-id feh --bg-fill ~/wallpaper.jpg &
exec --no-startup-id nm-applet &
exec --no-startup-id blueman-applet &
exec --no-startup-id udiskie --tray &
exec --no-startup-id /usr/lib/polkit-gnome/polkit-gnome-authentication-agent-1 &
exec --no-startup-id xss-lock -- i3lock -f -c 000000 &

### Status Bar ###
bar {
    status_command i3status
    # or if you installed polybar:
    # exec_always --no-startup-id polybar
}

### Screenshot ###
bindsym Print exec flameshot gui
bindsym Shift+Print exec flameshot full -p ~/Pictures/

### Volume ###
bindsym XF86AudioRaiseVolume exec pactl set-sink-volume @DEFAULT_SINK@ +5%
bindsym XF86AudioLowerVolume exec pactl set-sink-volume @DEFAULT_SINK@ -5%
bindsym XF86AudioMute exec pactl set-sink-mute @DEFAULT_SINK@ toggle

### Brightness ###
bindsym XF86MonBrightnessUp exec brightnessctl set +10%
bindsym XF86MonBrightnessDown exec brightnessctl set 10%-

### Lock ###
bindsym $mod+l exec i3lock -f -c 000000
\end{lstlisting}

\section{Sharing Config Between i3 and Sway (The Smart Way)}

The two configs are nearly identical --- you can maintain one shared config with separate override files for each WM.

\subsection{Structure}

\begin{lstlisting}
~/.config/i3/config             <- Shared config (keybindings, layouts, etc.)
~/.config/i3/conf.d/i3.conf    <- i3-only settings (polybar, picom, i3lock)
~/.config/sway/config           <- Symlink to the shared config
~/.config/sway/conf.d/sway.conf <- Sway-only settings (waybar, swaylock)
\end{lstlisting}

\subsection{Implementation}

\begin{lstlisting}
# Create the symlink
mkdir -p ~/.config/sway/conf.d ~/.config/i3/conf.d
ln -s ~/.config/i3/config ~/.config/sway/config
\end{lstlisting}

Add at the end of \texttt{\textasciitilde/.config/i3/config}:
\begin{lstlisting}
# Load WM-specific extras
include conf.d/*.conf
\end{lstlisting}

In \texttt{\textasciitilde/.config/i3/conf.d/i3.conf}:
\begin{lstlisting}
exec --no-startup-id picom -b &
exec --no-startup-id polybar &
exec --no-startup-id feh --bg-fill ~/wallpaper.jpg &
bindsym $mod+l exec i3lock -f -c 000000
\end{lstlisting}

In \texttt{\textasciitilde/.config/sway/conf.d/sway.conf}:
\begin{lstlisting}
exec waybar &
exec swaybg -i ~/wallpaper.jpg -m fill &
bindsym $mod+l exec swaylock -f -c 000000
output * bg ~/wallpaper.jpg fill
\end{lstlisting}

\section{Theming in i3}

i3 uses \texttt{lxappearance} instead of \texttt{nwg-look}:

\begin{lstlisting}
lxappearance
\end{lstlisting}

Same idea --- choose GTK theme + Icons + Cursor and click Apply. On X11 it works more straightforwardly without the \texttt{gsettings} complexity.

\section{picom (Compositor for i3)}

Without it you will not have transparency or visual effects. Minimal setup:

\begin{lstlisting}
# Launch from i3 config
exec --no-startup-id picom -b

# Or with a custom config
mkdir -p ~/.config/picom
nano ~/.config/picom/picom.conf
\end{lstlisting}

\begin{lstlisting}
# picom.conf -- minimal setup
backend = "glx";
vsync = true;
shadow = false;
fading = true;
fade-in-step = 0.05;
fade-out-step = 0.05;

# Transparency for terminal emulators
opacity-rule = [
    "90:class_g = 'tilix'",
    "90:class_g = 'Alacritty'"
];
\end{lstlisting}

\section{Summary: What to Install on Top of What You Have}

\begin{lstlisting}
# X11 base
sudo xbps-install xorg-minimal xinit

# i3 ecosystem
sudo xbps-install i3 polybar dunst i3lock feh picom flameshot \
  xclip arandr lxappearance xss-lock

# No need to reinstall:
# pipewire, nm-applet, blueman, udiskie, polkit-gnome, brightnessctl
\end{lstlisting}

\section{Quick Verification}

\begin{lstlisting}
# Verify Xorg
which startx
ls /usr/bin/Xorg

# Verify i3
which i3
i3 --version

# Launch
startx
\end{lstlisting}

% =============================================================
\end{document}
